\chapter{Development Environment}

\tikzset{pbox/.style={draw=cgray,fill=cgray!8,rounded corners=2pt,font=\sffamily\small,
    inner sep=4pt,align=center},
  parr/.style={-{Stealth[length=2mm]},thick,cgray}}
\providecommand{\tnode}[4]{\node[anchor=west,text=#4] at (#1,-#2*0.46) {\ttfamily\small #3};}

Writing code is half the job. The other half is the setup around it: where to get packages,
how to keep each project's packages separate, which editor to use, and how to save your
history safely. This chapter is the setup you need for every Python and \code{anthropic} project.

% ------------------------------------------------------------
\section{Find Python code}

Before writing something, check if it already exists. Look in this order:

\begin{center}
\begin{tabularx}{\linewidth}{@{}l >{\raggedright\arraybackslash}X >{\raggedright\arraybackslash}X@{}}
\toprule
\textbf{Source} & \textbf{What it is} & \textbf{Notes} \\
\midrule
Standard library & ships with Python, no install & start here (Chapter~12) \\
PyPI & the Python Package Index; what \code{pip} downloads from & more than half a million packages \\
GitHub & source code hosting & sort by stars; check the last commit date \\
Recipe collections & short snippets for one task & good for ideas, read before you paste \\
\bottomrule
\end{tabularx}
\end{center}

\begin{gamerule}[: judge a package before installing]
Open its PyPI page and check: recent releases, a source link, a reasonable download count, a
maintainer you can identify. Copy the name \emph{exactly}: fake packages with near-identical names
(\code{anthropic} versus a typo) exist to steal API keys.
\end{gamerule}

% ------------------------------------------------------------
\section{Install packages}

\begin{center}
%%MD-BEGIN
% ```mermaid
% flowchart LR
%     A["pip install anthropic"] --> B["PyPI<br/>(internet)"]
%     B --> C["download<br/>package + its dependencies"]
%     C --> D[".venv/Lib/site-packages"]
%     D --> E["import anthropic<br/>(found via sys.path)"]
% ```
%%MD-END
\begin{tikzpicture}[node distance=5mm]
  \node[pbox,draw=cblue,fill=cblue!10] (a) {\ttfamily\footnotesize pip install\\\ttfamily\footnotesize anthropic};
  \node[pbox,right=of a] (b) {PyPI\\{\footnotesize on the internet}};
  \node[pbox,right=of b] (c) {download package\\{\footnotesize + dependencies}};
  \node[pbox,right=of c,draw=cgreen,fill=cgreen!10] (d) {\ttfamily\footnotesize .venv\textbackslash Lib\\\ttfamily\footnotesize site-packages};
  \node[pbox,right=of d] (e) {\ttfamily\footnotesize import\\\ttfamily\footnotesize anthropic};
  \foreach \x/\y in {a/b,b/c,c/d,d/e} \draw[parr] (\x)--(\y);
\end{tikzpicture}
\end{center}

\begin{center}
\begin{tabularx}{\linewidth}{@{}l >{\raggedright\arraybackslash}X >{\raggedright\arraybackslash}X@{}}
\toprule
\textbf{Tool} & \textbf{What it does} & \textbf{When to pick it} \\
\midrule
\code{pip} & installs packages from PyPI; comes with Python & the default; learn it first \\
\code{venv} & standard-library module that makes virtual environments & every project \\
\code{virtualenv} & third-party ancestor of \code{venv} & old projects, extra features \\
Pipenv & pip + virtualenv + a lock file & legacy teams \\
Poetry & dependency management through \code{pyproject.toml} & publishing packages \\
Conda & installs Python \emph{and} non-Python libraries (science stack) & data science, GPU libraries \\
\code{uv} & very fast tool written in Rust; replaces pip, venv and more & new projects, if you want speed \\
\bottomrule
\end{tabularx}
\end{center}

\subsection{pip}

On Windows, prefer \code{python -m pip}. It runs pip \emph{for the exact interpreter} on the left, so
you cannot install into one Python and run another.

\begin{rawcode}
> python -m pip install anthropic              latest version
> python -m pip install "anthropic>=0.40"      minimum version
> python -m pip install pytest==8.3.4          exact version
> python -m pip install --upgrade anthropic    update
> python -m pip uninstall anthropic            remove
> python -m pip list                           what is installed here
> python -m pip freeze                         the same, in requirements format
\end{rawcode}

(The text on the right is a note, not part of the command. Not run here: these download from the internet and change an environment. The version numbers
are only examples of the syntax.)

\begin{trap}
Quote any specifier with \code{>}, \code{<} or \code{=} signs. In both PowerShell and cmd the shell
reads \code{>} as ``redirect output to a file'', so \code{pip install anthropic>=0.40} creates a file called
\code{=0.40}.
\end{trap}

\begin{center}
\begin{tabular}{@{}l l@{}}
\toprule
\textbf{Specifier} & \textbf{Meaning} \\
\midrule
\code{pkg} & any version, usually the newest \\
\code{pkg==1.4.2} & exactly this version \\
\code{pkg>=1.4} & this version or newer \\
\code{pkg>=1.4,<2} & newer than 1.4 but not 2.x \\
\code{pkg\textasciitilde{}=1.4.2} & compatible: at least 1.4.2, below 1.5 \\
\bottomrule
\end{tabular}
\end{center}

\textbf{Requirements file.} A text file listing the packages a project needs, one per line. Anyone
can rebuild the same setup with one command.

\begin{rawcode}
# requirements.txt
anthropic>=0.40
pytest>=8
python-dotenv

> python -m pip install -r requirements.txt
\end{rawcode}

\begin{gamerule}
Keep \code{requirements.txt} in the project and commit it to Git. Use \code{pip freeze > requirements.txt}
to record the \emph{exact} versions that work today, so a teammate (or you in six months) gets the same ones.
\end{gamerule}

\subsection{Other ways to install}

\begin{itemize}[leftmargin=*,itemsep=1pt]
\item \textbf{Native package managers} install programs, not just Python packages: Homebrew on
  macOS, \code{apt} on Linux, and on Windows the MSI installer or \code{winget}. Python itself
  on Windows usually comes this way, along with the \code{py} launcher.
\item \textbf{From source.} If a project is not on PyPI, download or \code{git clone} it, then run
  \code{python -m pip install .} inside its folder. (The book mentions \code{python setup.py install};
  that old command is no longer recommended.) You can also install straight from Git with
  \code{pip install git+https://...}.
\end{itemize}

\begin{trap}
Installing code means running someone else's code on your machine, with access to your files and
environment variables, including your \code{ANTHROPIC\_API\_KEY}. Pure-Python packages are readable, but
malware in them has happened. Install only what you trust.
\end{trap}

% ------------------------------------------------------------
\section{Virtual environments}

\subsection{Why you need one}

Project A may need \code{anthropic} 0.40, project B an older release. If both install into the one
system Python, they collide. A \textbf{virtual environment} is a private folder holding a Python
interpreter link plus its own \code{site-packages}. Each project gets its own.

\begin{center}
%%MD-BEGIN
% ```mermaid
% flowchart LR
%     subgraph sys["System Python 3.12 (shared, leave it clean)"]
%       L["python.exe + standard library"]
%     end
%     subgraph v1["project-a/.venv"]
%       P1["site-packages:<br/>anthropic 0.40, pytest"]
%     end
%     subgraph v2["project-b/.venv"]
%       P2["site-packages:<br/>older anthropic"]
%     end
%     v1 -. "home =" .-> sys
%     v2 -. "home =" .-> sys
% ```
%%MD-END
\begin{tikzpicture}[node distance=10mm]
  \node[pbox,text width=4.3cm] (sys) {\textbf{System Python 3.12}\\{\footnotesize \texttt{python.exe} and the standard library. Shared, so keep it clean.}};
  \node[pbox,draw=cgreen,fill=cgreen!10,text width=5.2cm,above right=0mm and 14mm of sys] (v1) {\textbf{project-a\textbackslash.venv}\\{\footnotesize \texttt{site-packages}: anthropic 0.40, pytest}};
  \node[pbox,draw=cgold,fill=cgold!10,text width=5.2cm,below right=0mm and 14mm of sys] (v2) {\textbf{project-b\textbackslash.venv}\\{\footnotesize \texttt{site-packages}: an older anthropic}};
  \draw[parr,dashed] (v1.west)--node[above,font=\scriptsize\sffamily,sloped]{\texttt{home =}}(sys.north east);
  \draw[parr,dashed] (v2.west)--(sys.south east);
\end{tikzpicture}
\end{center}

\subsection{Create and use one with \texttt{venv}}

\begin{rawcode}
> cd my-chat-app
> python -m venv .venv                  create the folder .venv
> .venv\Scripts\activate                cmd.exe
> .venv\Scripts\Activate.ps1            PowerShell  (see the trap below)
(.venv) > python -m pip install anthropic
(.venv) > deactivate                    leave the environment
\end{rawcode}

On macOS and Linux the activation command is \code{source .venv/bin/activate}, and the folder is
\code{bin} instead of \code{Scripts}. This is what the folder contains (real listing on Windows):

\begin{center}
%%MD-BEGIN
% ```text
% .venv/
% |-- pyvenv.cfg        <- records "home =" the base Python
% |-- Include/
% |-- Lib/site-packages <- packages for THIS project
% `-- Scripts/          <- python.exe, pip.exe, activate scripts
% ```
%%MD-END
\begin{tikzpicture}
  \tnode{0}{0}{.venv}{cgreen}
  \draw[cgray] (0.15,-0.3)--(0.15,-1.8);
  \foreach \i/\nm/\cmt in {1/pyvenv.cfg/records the base Python (home =),
      2/Include/C headers (rarely used),
      3/Lib\textbackslash site-packages/packages for THIS project,
      4/Scripts/python.exe\, pip.exe\, activate scripts}{
    \draw[cgray] (0.15,-\i*0.46)--(0.4,-\i*0.46);
    \tnode{0.45}{\i}{\nm}{ink}
    \node[font=\sffamily\footnotesize,text=cgray,anchor=west] at (5.2,-\i*0.46) {\cmt};}
\end{tikzpicture}
\end{center}

\subsection{What ``activate'' really does}

Your shell finds programs by searching the folders in the \code{PATH} environment variable, left to
right. Activation puts \code{.venv\textbackslash Scripts} at the \emph{front}, so typing \code{python} or \code{pip} finds the
project's copies first. Nothing else changes, and only in that terminal window.

\begin{center}
%%MD-BEGIN
% ```text
% before activate:  PATH = C:\Windows\System32 ; ...\Python312 ; ...
% after  activate:  PATH = ...\.venv\Scripts ; C:\Windows\System32 ; ...\Python312 ; ...
%                          ^ found first, so "python" means the venv's python
% ```
%%MD-END
\begin{tikzpicture}[node distance=2mm]
  \node[font=\sffamily\footnotesize,anchor=east] at (0,0) {before};
  \node[pbox,anchor=west] (b1) at (0.2,0) {\ttfamily\footnotesize C:\textbackslash Windows\textbackslash System32};
  \node[pbox,right=of b1] (b2) {\ttfamily\footnotesize ...\textbackslash Python312};
  \node[pbox,right=of b2] (b3) {\ttfamily\footnotesize ...};
  \node[font=\sffamily\footnotesize,anchor=east] at (0,-1.1) {after};
  \node[pbox,anchor=west,draw=cgreen,fill=cgreen!10] (a0) at (0.2,-1.1) {\ttfamily\footnotesize .venv\textbackslash Scripts};
  \node[pbox,right=of a0] (a1) {\ttfamily\footnotesize C:\textbackslash Windows\textbackslash System32};
  \node[pbox,right=of a1] (a2) {\ttfamily\footnotesize ...\textbackslash Python312};
  \node[font=\sffamily\scriptsize,text=cgreen!60!black,below=1mm of a0] {searched first};
\end{tikzpicture}
\end{center}

\begin{trap}
\textbf{PowerShell may refuse the script.} Windows can block \code{Activate.ps1} with an error about
``running scripts is disabled''. You do not need activation at all: call the interpreter directly,
\code{.venv\textbackslash Scripts\textbackslash python.exe -m pip install anthropic}. VS Code selects the interpreter for you
(command palette, \emph{Python: Select Interpreter}) and activates it in new terminals.
\end{trap}

\subsection{Am I inside the venv?}

Two attributes tell you. \code{sys.prefix} is the environment you run in; \code{sys.base\_prefix} is the Python
it was built from. They differ only inside a venv.

\begin{pycode}
# envcheck.py
import os
import sys

print(sys.executable.endswith("python.exe"))
print("in a venv:", sys.prefix != sys.base_prefix)

key = os.environ.get("DEMO_API_KEY")        # None if the variable is missing
print(type(key))
print("key is set" if key else "key is missing")

os.environ["DEMO_PORT"] = "8080"            # values are always str
port = os.environ["DEMO_PORT"]
print(type(port), int(port) + 1)

try:
    os.environ["NOPE_NOT_SET"]
except KeyError as err:
    print("KeyError:", err)
\end{pycode}
\begin{rawcode}
True
in a venv: True
<class 'NoneType'>
key is missing
<class 'str'> 8081
KeyError: 'NOPE_NOT_SET'
\end{rawcode}

(Run with the interpreter of a venv created by \code{python -m venv .venv}.)

\begin{typecheck}
\code{os.environ} behaves like a \code{dict[str, str]}: every value is a \textbf{string}, so convert with \code{int(...)}
yourself. \code{os.environ["X"]} raises \code{KeyError} when missing; \code{os.environ.get("X")} returns
\code{None}. That is why you check an \code{Optional[str]} before using an API key.
\end{typecheck}

\begin{gamerule}[: keep secrets out of code]
Put \code{ANTHROPIC\_API\_KEY} in an environment variable or a \code{.env} file, never in a \code{.py} file.
List \code{.env} and \code{.venv/} in \code{.gitignore} so neither is ever committed.
\end{gamerule}

\subsection{The other tools, in a sentence each}

\begin{itemize}[leftmargin=*,itemsep=1pt]
\item \textbf{virtualenv}: the older third-party tool that \code{venv} was modelled on; same idea.
\item \textbf{Pipenv}: runs pip and a virtualenv for you and writes a lock file of exact versions.
\item \textbf{Poetry}: manages dependencies and packaging from a \code{pyproject.toml} file; needs a recent Python 3.
\item \textbf{Conda}: a general package manager (not only Python), common for scientific stacks.
\item \textbf{uv}: a fast modern replacement for most of the above. Same ideas, different commands:
  \code{uv venv}, \code{uv pip install}.
\end{itemize}

% ------------------------------------------------------------
\section{Integrated development environments}

An IDE adds to a plain editor: autocomplete, inline docs, formatting and style checks, a file tree, source control,
test and debug buttons. Pick one and learn it well.

\begin{center}
\begin{tabularx}{\linewidth}{@{}l >{\raggedright\arraybackslash}X >{\raggedright\arraybackslash}X@{}}
\toprule
\textbf{Tool} & \textbf{Kind} & \textbf{Comment} \\
\midrule
Notepad++, vi/Vim, Neovim & text editors & fast and light; Vim keys take weeks to learn \\
Sublime Text, Zed & fast code editors & Zed has built-in AI assistant features \\
IDLE & Python's own small IDE & ships with Python; fine for first steps \\
PyCharm & full Python IDE & free Community edition; heavy but complete \\
\textbf{VS Code} & editor plus extensions & \textbf{free, most popular; you use it}; install the Python extension \\
Cursor & VS Code fork with AI built in & same look and settings as VS Code \\
IPython & improved interactive shell & \code{?} help, Tab completion, history lists \code{In} and \code{Out} \\
Jupyter, JupyterLab & notebooks: code, output and text in cells & data and experiments; also inside VS Code \\
\bottomrule
\end{tabularx}
\end{center}

\subsection{IPython and Jupyter}

IPython is a nicer version of the plain interactive Python prompt. Install it with \code{python -m pip install ipython}. Typing
\code{name?} shows the type, value and docstring of any object, and the \emph{Tab} key completes names.
Jupyter wraps this in a browser notebook (\code{pip install jupyterlab}, run \code{jupyter lab}):
you run small cells and see results below each one. Notebooks are excellent for trying out prompts and
looking at data.

\begin{trap}
A notebook keeps variables alive between cells, in the order you \emph{ran} them, not the order they appear.
A result can depend on a cell you deleted. Before trusting a notebook, restart the kernel and run all cells
from the top. For application code, use plain \code{.py} files.
\end{trap}

% ------------------------------------------------------------
\section{Source control}

Source control records every saved version of your project, so you can see what changed, undo a mistake
or work with others. \textbf{Git} is the standard; \textbf{Mercurial} (written in Python) is similar but
less common. Both are \emph{distributed}: every copy holds the full history. \textbf{GitHub} is a website
that hosts Git repositories.

\subsection{The three places a change lives}

\begin{center}
%%MD-BEGIN
% ```mermaid
% flowchart LR
%     W["working folder<br/>(your files)"] -- "git add" --> S["staging area<br/>(next snapshot)"]
%     S -- "git commit" --> R["repository<br/>(saved history)"]
%     W -. "git diff: what changed" .-> S
% ```
%%MD-END
\begin{tikzpicture}[node distance=22mm]
  \node[pbox,draw=cblue,fill=cblue!10,minimum width=3.2cm] (w) {working folder\\{\footnotesize the files you edit}};
  \node[pbox,right=of w,draw=cgold,fill=cgold!10,minimum width=3.2cm] (s) {staging area\\{\footnotesize the next snapshot}};
  \node[pbox,right=of s,draw=cgreen,fill=cgreen!10,minimum width=3.2cm] (r) {repository\\{\footnotesize saved history}};
  \draw[parr] (w)--node[above,font=\ttfamily\scriptsize]{git add}(s);
  \draw[parr] (s)--node[above,font=\ttfamily\scriptsize]{git commit}(r);
  \draw[parr,dashed] (w.south) to[bend right=25] node[below,font=\ttfamily\scriptsize]{git diff shows the gap} (s.south);
\end{tikzpicture}
\end{center}

\subsection{A first repository, step by step}

Real output from a scratch folder (Git 2.x on Windows):

\begin{rawcode}
> git init -b main
Initialized empty Git repository in .../repo/.git/

> echo print("Oops") > hello.py
> git add hello.py
> git status
On branch main
No commits yet
Changes to be committed:
        new file:   hello.py

> git commit -m "add hello script"
[main (root-commit) 8357896] add hello script
 1 file changed, 1 insertion(+)
 create mode 100644 hello.py
\end{rawcode}

Now edit \code{hello.py} (change \code{Oops} to \code{Hi}) and look at the change before saving it:

\begin{rawcode}
> git status
Changes not staged for commit:
        modified:   hello.py

> git diff
--- a/hello.py
+++ b/hello.py
@@ -1 +1 @@
-print("Oops")
+print("Hi")

> git add hello.py
> git commit -m "fix greeting"
[main 619bb0e] fix greeting
 1 file changed, 1 insertion(+), 1 deletion(-)

> git log --oneline
619bb0e fix greeting
8357896 add hello script
\end{rawcode}

(Git on Windows may also print \code{warning: LF will be replaced by CRLF}. It is about line endings
and is harmless here.)

\begin{center}
%%MD-BEGIN
% ```mermaid
% flowchart LR
%     C1["8357896<br/>add hello script"] --> C2["619bb0e<br/>fix greeting"] --> H["HEAD / main"]
% ```
%%MD-END
\begin{tikzpicture}[node distance=10mm]
  \node[pbox] (c1) {\ttfamily\footnotesize 8357896\\{\sffamily add hello script}};
  \node[pbox,right=of c1] (c2) {\ttfamily\footnotesize 619bb0e\\{\sffamily fix greeting}};
  \node[pbox,right=of c2,draw=cgreen,fill=cgreen!10] (h) {\ttfamily\footnotesize HEAD $\to$ main};
  \draw[parr] (c1)--node[above,font=\scriptsize\sffamily]{then}(c2);
  \draw[parr] (c2)--(h);
\end{tikzpicture}
\end{center}

\begin{center}
\begin{tabularx}{\linewidth}{@{}l X@{}}
\toprule
\textbf{Command} & \textbf{What it does} \\
\midrule
\code{git init} & turn the current folder into a repository \\
\code{git status} & show what is changed, staged, or untracked \\
\code{git add file} / \code{git add .} & put a file (or all changes) into the staging area \\
\code{git commit -m "msg"} & save the staged snapshot with a message \\
\code{git diff} & show line changes not yet staged \\
\code{git log} / \code{git log --oneline} & show history, newest first \\
\code{git clone url} & copy a whole repository from a host such as GitHub \\
\bottomrule
\end{tabularx}
\end{center}

\begin{trap}
\code{git commit} saves only what you staged. If you edit a file and commit without \code{git add}, Git
answers ``no changes added to commit'' and nothing is saved. Also: \code{git add .} stages \emph{everything},
including a \code{.env} file with a key. Write the \code{.gitignore} first:
\begin{rawcode}
.venv/
__pycache__/
.env
\end{rawcode}
\end{trap}

\begin{gamerule}
Commit small and often, with a message that says \emph{what} changed (``fix greeting''), not ``stuff''.
Commit when the tests pass. Review with \code{git diff} before every commit.
\end{gamerule}

\begin{boss}[: ``which Python am I running?'']
Most setup bugs are a mismatch: you installed a package into one Python and run your code with another
(\code{ModuleNotFoundError: No module named 'anthropic'} right after a successful install). Ask three questions:
(1) \code{import sys; print(sys.executable)}: which interpreter is running? (2) \code{sys.prefix != sys.base\_prefix}: is it a venv?
(3) \code{python -m pip --version}: does pip belong to the same interpreter? Always install with
\code{python -m pip} and make VS Code use the \code{.venv} interpreter, and the mismatch disappears.
\end{boss}

\begin{practice}
\begin{enumerate}[leftmargin=*,itemsep=2pt]
\item Create a project folder with a \code{.venv}. Which two commands activate it in PowerShell and in cmd, and how do you leave it?
\item You ran \code{pip install anthropic} but Python says the module is missing. List three causes you would check.
\item Write a \code{requirements.txt} that needs any \code{anthropic}, \code{pytest} 8 or newer, and nothing else. How does a teammate install it?
\item What does \code{git add} do that \code{git commit} does not? Why is there a staging area at all?
\item Write the lines for a \code{.gitignore} in an LLM project and say why each line is there.
\item Read an API key from an environment variable. What is the type of the result when the variable is missing, and what do you do about it?
\end{enumerate}
\end{practice}

\begin{achievement}{Environment Engineer}
\begin{itemize}[leftmargin=*,itemsep=1pt]
\item Create a virtual environment, install packages into it, and prove which Python you are using.
\item Record and rebuild a project's dependencies with \code{requirements.txt}.
\item Keep secrets out of code and out of Git with environment variables and \code{.gitignore}.
\item Make, review and list commits with Git.
\end{itemize}
\end{achievement}

%%MD-BEGIN
% **Related:** [[12 - Modules and Packages]] (what pip installs), [[14 - Type Hints and Documentation]] (mypy runs in the venv), [[15 - Testing]] (pytest, CI), [[16 - Debugging]] (find which Python), [[20 - Files]] (reading `.env`).
%%MD-END
